\scanpage{134}
\mspage{-III- \handwritten{3} -}
\begin{flushright}n°\handwritten{383}\end{flushright}
pour $G^{\unsure{o|0|\circ}}$ est le plan $P$ engendré par $e_i$ et $e_j$. Les hyperplans de la forme $g(H_i)$ et $g(H_j)$, pour $g\in G^{\unsure{o|0|\circ}}$, contiennent tous le sous-espace $P'$ orthogonal à $P$, et forment une partie $\underline{H}^{\unsure{o|0|\circ}}$ de $\underline{H}$ qui jouit, vis à vis de $G^{\unsure{o|0|\circ}}$, des propriétés 1) et 2) ci-dessus. Comme $\underline{H}^{\unsure{o|0|\circ}}\handwritten{\subset}\underline{H}$, il existe une chambre $C^{\unsure{o|0|\circ}}$ définie par $\underline{H}^{\unsure{o|0|\circ}}$ et contenant $C$ ; il est clair que $C^{\unsure{o|0|\circ}}$ est un des dièdres ayant $H_i$ et $H_j$ pour murs. Ainsi $C^{\unsure{o|0|\circ}}\handwritten{\cap}P$ est un secteur angulaire, limité par deux demi-droites $D_i\handwritten{\subset}H_i$ et $D_j\handwritten{\subset}H_j$. Posant $\handwritten{\vartheta}=(\handwritten{\widehat}{D_i,D_j})$, $s_j s_i$ est la rotation de $2\handwritten{\vartheta}$ autour de $P'$, et les transformées de $C^{\unsure{o|0|\circ}}$ par $G^{\unsure{o|0|\circ}}$ sont ses transformées par les rotations d'angle $n\handwritten{\vartheta}$ $(n\in Z)$ ; comme deux transformées distinctes de $C^{\unsure{o|0|\circ}}$ sont disjointes, il résulte du \S\ 1, n°3, lemme 9 que $\handwritten{\vartheta}$ est de la forme $\handwritten{\pi}/m$ $(m\in Z)$ ; on a nécessairement $m=m_{ij}$. Par définition de $e_i$ et $e_j$, et avec une orientation convenable de $P$, on a $(\handwritten{\widehat}{D_i,e_i})=-\handwritten{\pi}/2$, et $(\handwritten{\widehat}{D_j,e_j})=+\handwritten{\pi}/2$. D'où, par la formule de Chasles (Alg.IX, \S\ 10, n°3, (9)), $(\handwritten{\widehat}{e_i,e_j})=(\handwritten{\widehat}{e_i,D_i})+(\handwritten{\widehat}{D_i,D_j})+(\handwritten{\widehat}{D_j,e_j})=\handwritten{\pi}+(\handwritten{\widehat}{D_i,D_j})$. On en déduit $(e_i\handwritten{|}e_j)=\cos(\handwritten{\widehat}{e_i,e_j})=-\cos(\handwritten{\widehat}{D_i,D_j})=-\cos(\handwritten{\pi}/m_{ij})$. La seconde assertion résulte de la monotonie de la fonction cosinus dans $[0,\handwritten{\pi}/2]$ et de la formule $\cos(\handwritten{\pi}/3)=1/2$ (référence introuvable !!).

PROPOSITION 3.- \uline{Les vecteurs} $e_1,\ldots,e_m$ \uline{forment une base du sous-espace essentiel} $W$ \uline{de} $E$.

En effet l'ensemble $C'$ formé par $0$ et par les $y\in E$ tels que $(y\handwritten{|}x)>0$ pour tout $x\in C$ est un cône convexe pointé, et on a $C'\handwritten{\cap}(-C')=\handwritten{\{}0\handwritten{\}}$ ; donc $C'$ est l'ensemble des éléments positifs pour une relation d'ordre faisant de $E$ un espace vectoriel ordonné (E.V.T., II, \S\ 1, n°5). D'après (2) on a $e_i>0$ pour tout $i$, et d'après la prop.2, on a $(e_i\handwritten{|}e_j)\handwritten{\leq}0$ pour $i\handwritten{\ne}j$. Le lemme 6 du \S\ 1, n°2 montre alors que les $e_i$ sont linéairement indépendants. D'autre part il est clair que les $e_i$ appartiennent à $W$. Enfin si un élément de $E$ est orthogonal à $e_1,\ldots,e_m$, il est fixe pour les
% archive sheet number: 133
